\documentclass[a4paper, 11pt]{article}
\usepackage{geometry}
\geometry{letterpaper, margin=1in}
\usepackage{amsmath}
\usepackage{amssymb}  
\usepackage{amsthm}
\usepackage{ulem} 
\usepackage{graphicx}
\usepackage{enumitem} % use for making lettered list 
\usepackage{bbm} % use for making the 1 identity operator EX: \mathbbm{1}
\usepackage{subfig}
\graphicspath{ {images/} }

% format to allow bolded theorems, corollaries, etc... 
\newtheorem*{theorem}{Theorem}
\newtheorem*{corollary}{Corollary}
\newtheorem*{lemma}{Lemma}
\newtheorem*{definition}{Definition}
\newtheorem*{Example}{Example} 
\newtheorem*{Remark}{Remark}

% stop typing \mathbb a thousand times 
\newcommand{\R}{\mathbb{R}}
\newcommand{\C}{\mathbb{C}}
\newcommand{\F}{\mathbb{F}}
\newcommand{\Mat}[2]{\mathcal{M}_{#1\times#2}}
\newcommand{\proj}[2]{\operatorname{proj}_{#2}(#1)}

% change margins for solution
\newenvironment{solution}{%
	\begin{list}{}{%
			\setlength{\topsep}{0pt}%
			\setlength{\leftmargin}{1.5cm}%
			\setlength{\rightmargin}{1.5cm}%
			\setlength{\listparindent}{\parindent}%
			\setlength{\itemindent}{\parindent}%
			\setlength{\parsep}{\parskip}%
		}%
		\item[]}{\end{list}}

\begin{document}
%Header-Make sure you update this information!!!!
\noindent
\large\textbf{Homework 8} \hfill \textbf{John Waczak} \\
\normalsize MTH 443 \hfill  Date: \today \\
Dr. Schmidt  
\par\noindent\rule{\textwidth}{0.4pt} \\\\

\noindent 1.) Let $V$ be the real inner product space given by $\mathcal{P}_2(\R)$ with the inner product
\begin{equation*}
  \langle f, g \rangle = \int\limits_0^1 f(x)g(x)dx
\end{equation*}
Perform Gram-Schmidt orthonormalization on the set $\mathcal{S}=\left\{ x^2-1, x^2+3, x-2 \right\}$ \\

\begin{solution}
 \noindent For ease of calculation, let us define the elements of $\mathcal{S}$ to be
  \begin{align*}
    f_1(x) &= x^2-1\\
    f_2(x) &= x^2+3 \\
    f_3(x) &= x-2
  \end{align*}
  The functions in the desired ONB will then be defined as $\{g_1(x), g_2(x), g_3(x)\}$. Recall that the definition of the projection of the vector v onto the vector u is given as
  \begin{equation*}
    \proj{u}{v} = \frac{\langle v,u \rangle}{\langle u, u \rangle}u
  \end{equation*}
  Now to perform Gram-Schmidt orthonormalization, we first take $\mathcal{S}$ and orthogonalize it. We accomplish this by first defining $f_1(x)=g_1(x)$. Then we construct the rest of the orthogonalized vectors $g_2, g_3$ by utilizing the projection defined above. Pictorially this is
  \begin{figure}[!hbt]
    \centering
    \includegraphics[width=0.3\textwidth]{gramSchmidt.pdf}
    \caption{The Gram-Schmidt orthogonalization process}
    \label{fig:gramSchmidt}
  \end{figure}
  For each subsequent vector, we just subtract of more projections so that in the end, we have
  \begin{align*}
    g_1(x) &= f_1(x) \\
    g_2(x) &= f_2(x)-\proj{f_2}{g_1}\\
    g_3(x) &= f_3(x)-\proj{f_3}{g_1}-\proj{f_3}{g_2}
  \end{align*}

  \noindent Given our functions from $\mathcal{S}$, this becomes:
  \begin{align*}
    g_1(x) &= x^2-1\\
    g_2(x) &= x^2+3 - \frac{\int_0^1(x^2+3)(x^2-1)dx}{\int_0^1(x^2-1)^2dx}(x^2-1) \\
    &= x^2+3 - \frac{\frac{-32}{15}}{\frac{8}{15}}(x^2-1) \\
    &= x^2+3+4(x^2-1) \\
    &= 5x^2-1 \\
    g_3(x) &= x-2 - \frac{\int_0^1(x-2)(x^2-1)dx}{\int_0^1(x^2-1)^2dx}(x^2-1) - \frac{\int_0^1(x-2)(5x^2-1)dx}{\int_0^1(5x^2-1)^2dx}(5x^2-1)\\
    &= x-2 - \frac{\frac{13}{12}}{\frac{8}{15}}(x^2-1) -\frac{\frac{-7}{12}}{\frac{8}{3}}(5x^2-1)\\
    &= -\frac{3}{16}+x-\frac{15}{16}x^2 
  \end{align*}

  \noindent To ensure we haven't made a mistake it is valuable to check that these new vectors are in fact linearly independent. Writing the matrix formed by the coordinate representation of each vector in $\R^3$ we have
  \begin{equation*}
    \det\begin{pmatrix}
    -1 & -1 & -\frac{3}{16} \\
    0 & 0 & 1 \\
    1 & 5 & -\frac{15}{16}
    \end{pmatrix} = -1(5)+1(-1) = 4 \neq 0
  \end{equation*}
  Therefore, we can be confident that we do in fact still have a basis. Now that we have orhtogonalized, we need to normalize it. To do this we calculate the lengths of each vector.
  \begin{align*}
    ||g_1|| &= \sqrt{\langle g_1, g_1 \rangle} \\
    &= \left( \int_0^1(x^2-1)^2dx\right)^{\frac{1}{2}} = 2\sqrt{\frac{2}{15}}\\
    ||g_2|| &= \left(\int_0^1(5x^2-1)^2dx \right)^{\frac{1}{2}} = 2\sqrt{\frac{2}{3}}\\
    ||g_3|| &= \left( \int_0^1\left( -\frac{3}{16}+x-\frac{15}{16}x^2   \right)^2dx  \right)^{\frac{1}{2}} = \frac{1}{8\sqrt{3}} 
  \end{align*}
  Therefore, our orthonormalized $\mathcal{S}$ is
  \begin{equation*}
    \mathcal{S}_\text{ONB} = \left\{ \frac{g_1}{||g_1||}, \frac{g_2}{||g_2||},  \frac{g_3}{||g_3||}  \right\} = \left\{ \frac{1}{2}\sqrt{\frac{15}{2}}(x^2-1), \frac{1}{2}\sqrt{\frac{3}{2}}(5x^2-1), -\frac{1}{2}\sqrt{3}(3-16x+15x^2)\right\}
  \end{equation*}
\end{solution}


\noindent 2.) Consider $\mathcal{M}_2(\C)$ with the inner product
\begin{equation*}
  \langle A, B \rangle = \text{tr}B^\star A,
\end{equation*}
where $\text{tr} M = \sum_i^n M_ii$ and $M^\star = \bar{M}^t$ (conjugate transpose). Find the `closest` element of the complex symmetric to $A=\begin{pmatrix} 1 & -i \\ i & 1 \end{pmatrix}$. \\

\begin{solution}
  \noindent Recall that a matrix $M$ is symmetric if $M=M^t$. We can find the 'closest' matrix to A by finding a ONB for the subspace of symmetric 2x2 complex matrices and then finding expansion of $A$ in terms of this basis. The result will be a symmetric matrix that is not equal to A but is the closest we can get while still being symmetric. Therefore, we will first find an ONB for the symmetric 2x2 matrices.\\

  \noindent From the above definition it is easy to see that a general symmetric 2x2 matrix may be written as
  \begin{equation*}
    M = \begin{pmatrix} a & b \\ b & c \end{pmatrix}
  \end{equation*}
  with $a,b,c\in\C$. Thus, this subspace is defined by three free parameters and therefore a suitable basis is
  \begin{equation*}
    \mathcal{B} = \left\{ M_1=\begin{pmatrix} 1 & 0 \\ 0 & 0 \end{pmatrix}, M_2=\begin{pmatrix} 0 & 1 \\ 1 & 0 \end{pmatrix}, M_3=\begin{pmatrix}0 & 0 \\ 0 & 1 \end{pmatrix}   \right\}
  \end{equation*}

  \noindent Before we can try and expand A in this basis, we must verify that it is orthonormal under our inner product $\langle \cdot, \cdot \rangle$.
  \begin{align*}
    \langle M_1, M_2 \rangle &= \text{tr}\begin{pmatrix}0 & 1 \\ 1 & 0\end{pmatrix}\begin{pmatrix}1 & 0 \\ 0 & 0 \end{pmatrix} \\
      &=\text{tr}\begin{pmatrix} 0 & 0 \\ 1 & 0 \end{pmatrix} \\
      &= 0 
  \end{align*}
  \begin{align*}
    \langle M_3, M_2 \rangle &= \text{tr}\begin{pmatrix}0 & 1 \\ 1 & 0\end{pmatrix}\begin{pmatrix}0 & 0 \\ 0 & 1 \end{pmatrix} \\
      &=\text{tr}\begin{pmatrix} 0 & 1 \\ 0 & 0 \end{pmatrix} \\
      &= 0 
  \end{align*}
  \begin{align*}
    \langle M_3, M_1 \rangle &= \text{tr}\begin{pmatrix}1 & 0 \\ 0 & 0\end{pmatrix}\begin{pmatrix}0 & 0 \\ 0 & 1 \end{pmatrix} \\
      &=\text{tr}\begin{pmatrix} 0 & 0 \\ 0 & 0 \end{pmatrix} \\
      &= 0 
  \end{align*}
\begin{align*}
    \langle M_1, M_1 \rangle &= \text{tr}\begin{pmatrix}1 & 0 \\ 0 & 0\end{pmatrix}\begin{pmatrix}1 & 0 \\ 0 & 0 \end{pmatrix} \\
      &=\text{tr}\begin{pmatrix} 1 & 0 \\ 0 & 0 \end{pmatrix} \\
      &= 1 
  \end{align*}
\begin{align*}
    \langle M_2, M_2 \rangle &= \text{tr}\begin{pmatrix}0 & 1 \\ 1 & 0\end{pmatrix}\begin{pmatrix}0 & 1 \\ 1 & 0 \end{pmatrix} \\
      &=\text{tr}\begin{pmatrix} 1 & 0 \\ 0 & 1 \end{pmatrix} \\
      &= 2 
  \end{align*}
\begin{align*}
    \langle M_3, M_3 \rangle &= \text{tr}\begin{pmatrix}0 & 0 \\ 0 & 1\end{pmatrix}\begin{pmatrix}0 & 0 \\ 0 & 1 \end{pmatrix} \\
      &=\text{tr}\begin{pmatrix} 0 & 0 \\ 0 & 1 \end{pmatrix} \\
      &= 1 
  \end{align*}

\noindent Thus we an see $\mathcal{B}$ is almost an ONB except for $M_2$ therefore, dividing each matrix by its norm gives the desired ONB
\begin{equation*}
  \mathcal{B}_\text{ON} = \left\{ m_1=\begin{pmatrix} 1 & 0 \\ 0 & 0 \end{pmatrix}, m_2=\begin{pmatrix} 0 & \frac{1}{\sqrt{2}} \\ \frac{1}{\sqrt{2}} & 0 \end{pmatrix}, m_3=\begin{pmatrix}0 & 0 \\ 0 & 1 \end{pmatrix}   \right\}
\end{equation*}

\noindent Now with this ONB, the 'clostest' matrix to $A$ which we will call $\tilde{A}$ is given by
\begin{align*}
  \tilde{A} &= \langle A, m_1\rangle m_1 + \langle A, m_2 \rangle m_2 + \langle A, m_3 \rangle m_3\\
  \langle A, m_1 \rangle &= \text{tr}\begin{pmatrix}1 & 0 \\ 0 & 0 \end{pmatrix}\begin{pmatrix}1 & -i \\ i & 1\end{pmatrix} \\
    &= \text{tr}\begin{pmatrix}1 & -i \\ 0 & 0 \end{pmatrix}\\
    &= 1\\
   \langle A, m_2 \rangle &= \text{tr}\begin{pmatrix}0 & 1 \\ 1 & 0 \end{pmatrix}\begin{pmatrix}1 & -i \\ i & 1\end{pmatrix} \\
    &= \text{tr}\begin{pmatrix}i & 1 \\ 1 & -i \end{pmatrix}\\
    &= i-i = 0 \\
    \langle A, m_3 \rangle &= \text{tr}\begin{pmatrix}0 & 0 \\ 0 & 1 \end{pmatrix}\begin{pmatrix}1 & -i \\ i & 1\end{pmatrix} \\
    &= \text{tr}\begin{pmatrix}0 & 0 \\ i & 1 \end{pmatrix}\\
    &= 1 \\
    \Rightarrow \tilde{A} &= \begin{pmatrix}1 & 0 \\ 0 & 0 \end{pmatrix} + \begin{pmatrix}0 & 0 \\ 0 & 1 \end{pmatrix} = \begin{pmatrix}1 & 0 \\ 0 & 1 \end{pmatrix} 
\end{align*}

\noindent Thus, we have found the closest symmetric 2x2 $\C$-valued matrix to $A$. \\
\end{solution}

\noindent 3.) Let $V$ be the real inner product space given by $\mathcal{P}_2(\R)$ with the inner product
\begin{equation*}
  \langle f,g \rangle = \int\limits_{-1}^1 f(x)g(x)dx
\end{equation*}
Let $\ell:V\to\R$ by $f\mapsto f'(0)+f(1)$. Find $g\in$ V, we have $\ell(f)=\langle f,g \rangle$\\

\begin{solution}
  \noindent Let $f,g \in V$ such that
  \begin{align*}
    f(x) &= a_0 + a_1x + a_2x^2 \\
    g(x) &= b_0 + b_1x + b_2x^2
  \end{align*}
  Differentiating $f$ and substituting into the equation of the problem statement yields,
  \begin{equation*}
    a_1 + a_0 + a_1 + a_2 = a_0+2a_1+a_2
  \end{equation*}
  The right hand side of the equation becomes
  \begin{align*}
    \langle f,g \rangle &= \int\limits_{-1}^1 (a_0+a_1x+a_2x^2)(b_0+b_1x+b_2x^2)dx \\
    &=\int\limits_{-1}^1 a_0b_0 + a_1b_0x + a_2b_0x^2 + a_0b_1x + a_1b_1x^2+a_2b_1x^3+a_0b_2x^2+a_1b_2x^3+a_2b_2x^4 dx\\
    &=2a_0b_0 + \frac{2a_2b_0}{3} + \frac{2a_1b_1}{3} + \frac{2a_0b_2}{3} + \frac{2a_2b_2}{5} \\
    &= \left(2b_0+\frac{2}{3}b_2\right)a_0 + \left(2\frac{b_1}{3}\right)a_1 + \left(2\frac{b_0}{3}+2\frac{b_2}{5}   \right)a_2
  \end{align*}
  Now in order for the two sides to be equal, we need the coefficients in front of each $a_i$ to be equal. Therefore, we have the following system
  \begin{equation*}
    \begin{cases}
      2b_0+2\frac{b_2}{3} = 1 \\\\
      2\frac{b_1}{3} = 2 \\\\
      2\frac{b_0}{3}+2\frac{b_2}{5}
      \end{cases}
  \end{equation*}
  This leads to
  \begin{align*}
    b_1 &= 3 \\
    2b_0 + \frac{2}{3}b_2 &= \frac{2}{3}b_0 + \frac{2}{5}b_2\\
    \frac{4}{3}b_0 &= \left(\frac{6}{15}-\frac{10}{15}\right)b_2 \\
    b_0 &= -\frac{1}{5}b_2\\
    &\vdots
  \end{align*}
  \begin{align*}
    &  \vdots \\
    1 &= \frac{2}{3}\left(-\frac{1}{5}b_2\right)+\frac{2}{5}b_2 \\
    b_2 &= \frac{1}{\frac{6}{15}-\frac{2}{15}} \\
    &= \frac{15}{4} \\
    b_0 &= -\frac{1}{5}\frac{15}{4} \\
    &= -\frac{3}{4}
  \end{align*}
  Thus, we have identified all the $b_i$'s and have therefore found the form for $g(x)$. We have
  \begin{equation*}
    g(x) = -\frac{3}{4} + 3x + \frac{15}{4}x^2
  \end{equation*}

\end{solution}


\end{document}



















